% Manual do CSG. Compile com `npm run manual` (docs/build-manual.mjs), que:
%   - avalia todo bloco \begin{csg}[run] com o interpretador e o núcleo reais;
%   - gera gen/functions.tex e gen/symbols.tex a partir do código;
%   - renderiza as figuras (\csgfig) com o visualizador da aplicação;
%   - colore os blocos csg com o mesmo realce do editor e compila com LuaLaTeX.
% O arquivo também compila diretamente (lualatex manual.tex), com código sem cores.
\documentclass[11pt,a4paper]{article}
\usepackage[a4paper,top=24mm,bottom=24mm,left=22mm,right=22mm,headheight=14pt]{geometry}
\usepackage[brazilian,shorthands=off]{babel}

% ---------- fontes ----------
\usepackage{fontspec}
\usepackage{unicode-math}
% símbolos que o Libertinus Serif não tem caem no Libertinus Math e na JuliaMono
\directlua{luaotfload.add_fallback("csgfallback",
  {"LibertinusMath-Regular.otf:mode=node;", "JuliaMono-Regular.ttf:mode=node;"})}
\setmainfont{LibertinusSerif}[Extension=.otf, UprightFont=*-Regular, BoldFont=*-Semibold,
  ItalicFont=*-Italic, BoldItalicFont=*-SemiboldItalic, RawFeature={fallback=csgfallback}]
\setsansfont{LibertinusSans}[Extension=.otf, UprightFont=*-Regular, BoldFont=*-Bold,
  ItalicFont=*-Italic, RawFeature={fallback=csgfallback}]
\setmathfont{LibertinusMath-Regular.otf}
\setmonofont{JuliaMono}[Path=fonts/, Extension=.ttf, UprightFont=*-Regular, BoldFont=*-SemiBold,
  ItalicFont=*-RegularItalic, BoldItalicFont=*-SemiBoldItalic, Scale=MatchLowercase,
  Contextuals=NoAlternate, Ligatures={NoCommon,NoRequired}]

% ---------- pacotes ----------
\usepackage{xcolor}
\usepackage{graphicx}
\usepackage{fvextra}
\usepackage[most]{tcolorbox}
\usepackage{booktabs,array,longtable}
\usepackage{enumitem}
\usepackage{titlesec}
\usepackage{fancyhdr}
\usepackage[font=small,labelfont={bf,sf},labelsep=endash]{caption}
\usepackage{microtype}
\usepackage[hidelinks]{hyperref}

% ---------- cores (as do tema claro do editor) ----------
\definecolor{accent}{HTML}{2F5FB3}
\definecolor{muted}{HTML}{666666}
\definecolor{codebg}{HTML}{F4F5F7}
\definecolor{rulec}{HTML}{D8DBE0}
\definecolor{notebg}{HTML}{FDF8E6}
\definecolor{noteframe}{HTML}{E0C36A}
\definecolor{tkcomment}{HTML}{8A8A85}
\definecolor{tknumber}{HTML}{A5521A}
\definecolor{tkkeyword}{HTML}{8B2FA8}
\definecolor{tkdef}{HTML}{1F5FBF}
\definecolor{tkbuiltin}{HTML}{0F7A6C}
\definecolor{tkset}{HTML}{C4372F}
\definecolor{tkiso}{HTML}{B26A00}
\definecolor{tkbracket}{HTML}{6B6B70}
\hypersetup{colorlinks, linkcolor=accent, urlcolor=accent, citecolor=accent,
  pdftitle={CSG — Manual da linguagem e da aplicação}, pdfsubject={Modelagem de sólidos por CSG}}

% ---------- página ----------
\setlength{\parindent}{0pt}
\setlength{\parskip}{5pt plus 1pt}
\setlist{itemsep=1pt, topsep=3pt, parsep=1pt, leftmargin=16pt}
\pagestyle{fancy}
\fancyhf{}
\fancyhead[R]{\sffamily\footnotesize\color{muted} CSG · manual}
\fancyfoot[C]{\small\color{muted}\thepage}
\renewcommand{\headrulewidth}{0pt}

\titleformat{\section}{\sffamily\bfseries\LARGE\color{accent}}{\thesection}{0.55em}{}[{\color{accent}\vspace{2pt}\titlerule[1.2pt]}]
\titleformat{\subsection}{\sffamily\bfseries\large}{\thesubsection}{0.5em}{}
\titleformat{\subsubsection}{\sffamily\bfseries\normalsize\color{black!80}}{}{0pt}{}
\titlespacing*{\section}{0pt}{0pt}{10pt}
\titlespacing*{\subsection}{0pt}{14pt}{4pt}
\titlespacing*{\subsubsection}{0pt}{10pt}{3pt}
\AddToHook{cmd/section/before}{\clearpage}

% ---------- código ----------
% inline: \code{...} lê o argumento literalmente (\, #, { } etc. são permitidos, se as chaves forem balanceadas)
\NewDocumentCommand\code{v}{\texttt{#1}}
\newcommand\kbd[1]{{\setlength{\fboxsep}{1.2pt}\setlength{\fboxrule}{0.4pt}%
  \raisebox{0.5pt}{\fcolorbox{black!35}{black!2}{\sffamily\footnotesize\vphantom{Ay}#1}}}}
\newtcolorbox{csgbox}{enhanced, colback=codebg, colframe=accent, boxrule=0pt, leftrule=2.5pt,
  arc=0pt, outer arc=0pt, left=7pt, right=6pt, top=2pt, bottom=2pt, before skip=5pt, after skip=8pt}
% blocos com realce (gerados pelo build) e sem realce
\DefineVerbatimEnvironment{csghl}{Verbatim}{commandchars=\\\{\}, fontsize=\small, breaklines, breaksymbolleft={}}
\DefineVerbatimEnvironment{csgpl}{Verbatim}{fontsize=\small, breaklines, breaksymbolleft={}}
% fonte: \begin{csg}[run, id=nome] … \end{csg}  (run = o build avalia o programa)
\makeatletter
\define@key{FV}{run}[]{}\define@key{FV}{id}{}\define@key{FV}{plain}[]{}
\makeatother
\DefineVerbatimEnvironment{csg}{Verbatim}{fontsize=\small, breaklines, frame=leftline,
  framerule=2.5pt, rulecolor=\color{accent}, xleftmargin=4pt}
\newcommand\tkBS{\char92\relax}
\newcommand\tkLB{\char123\relax}
\newcommand\tkRB{\char125\relax}
\newcommand\tkcomment[1]{\textcolor{tkcomment}{\itshape #1}}
\newcommand\tknumber[1]{\textcolor{tknumber}{#1}}
\newcommand\tkkeyword[1]{\textcolor{tkkeyword}{\bfseries #1}}
\newcommand\tkdef[1]{\textcolor{tkdef}{\bfseries #1}}
\newcommand\tkbuiltin[1]{\textcolor{tkbuiltin}{#1}}
\newcommand\tkset[1]{\textcolor{tkset}{\bfseries #1}}
\newcommand\tkiso[1]{\textcolor{tkiso}{\bfseries #1}}
\newcommand\tklogic[1]{\textcolor{tkkeyword}{#1}}
\newcommand\tkmath[1]{\textcolor{tknumber}{#1}}
\newcommand\tkbracket[1]{\textcolor{tkbracket}{#1}}
\newcommand\tklatex[1]{\textcolor{tkcomment}{#1}}
\newcommand\tkident[1]{#1}
\newcommand\tkop[1]{#1}

% ---------- tabelas, notas, figuras ----------
\newcolumntype{L}[1]{>{\raggedright\arraybackslash}p{\dimexpr#1\linewidth-\tabcolsep\relax}}
\setlength{\LTleft}{0pt}\setlength{\LTright}{0pt}\setlength{\LTpre}{4pt}\setlength{\LTpost}{6pt}
\newcommand\colhead[1]{{\sffamily\footnotesize\color{muted}\MakeUppercase{#1}}}
\renewcommand{\arraystretch}{1.22}
\newtcolorbox{nota}{colback=notebg, colframe=noteframe, boxrule=0.6pt, arc=1pt,
  left=7pt, right=7pt, top=4pt, bottom=4pt, before skip=6pt, after skip=10pt, fontupper=\small}
% \csgfig{parâmetros da renderização}{arquivo em img/ sem extensão}{fração da largura}{legenda}
% parâmetros: pre=ID (programa do bloco com id=ID), size=LxA, dist=, view=iso|top|front|side
\newcommand\csgfig[4]{\begin{figure}[!ht]\centering
  \includegraphics[width=#3\linewidth,height=80mm,keepaspectratio]{img/#2.png}
  \caption{#4}\end{figure}}

\begin{document}

% ======================================================= capa
\begin{titlepage}
\vspace*{38mm}
{\fontsize{64}{64}\selectfont\ttfamily\color{accent} ⋃ ∩ ∖\par}
\vspace{16pt}
{\sffamily\bfseries\fontsize{34}{38}\selectfont CSG\par}
\vspace{8pt}
{\sffamily\Large\color{muted} Modelagem de sólidos com uma linguagem de inspiração matemática\par}
\vspace{2pt}
{\sffamily\large\color{muted} Manual da linguagem e da aplicação\par}
\vspace{28pt}
\begin{minipage}{130mm}
\begin{csg}[run]
furos ≔ ⋃{ Rz(k · 45°) · T(24, 0, 0) · cylinder(2, 10) | k ∈ 0..<8 }
peça  ≔ cylinder(30, 4) ∖ furos ∖ sphere(8)
⊢ peça
\end{csg}
\end{minipage}
\vfill
{\sffamily\small\color{muted} Versão de \today\par}
\end{titlepage}

{\hypersetup{linkcolor=black}\tableofcontents}

% =======================================================
\section{Introdução}\label{sec:intro}

\textbf{CSG} (\emph{Constructive Solid Geometry}, geometria construtiva de sólidos) descreve um objeto tridimensional como
uma expressão: sólidos simples, como cubos, esferas e cilindros, são posicionados por transformações e combinados por
operações da teoria dos conjuntos: união $\cup$, interseção $\cap$ e diferença $\setminus$.
Esta aplicação leva essa ideia a sério. O modelo é escrito numa linguagem cuja notação é a da matemática, e cada
definição é uma igualdade entre conjuntos de pontos do espaço.

Um programa curto já produz uma peça imprimível:

\begin{csg}[run, id=flange-mini]
r ≔ 30                                       # raio da base
furo(θ) ≔ Rz(θ) ∘ T(r - 6, 0, 0) · cylinder(r = 2, h = 10)
base    ≔ cylinder(r, 4)
furos   ≔ ⋃{ furo(k · 45°) | k ∈ 0..<8 }
pilares ≔ ⋃{ T(⟨20sx, 20sy, 0⟩) · box(4, 4, 20) | sx, sy ∈ {-1, 1} }
peça    ≔ (base ∖ furos) ∪ pilares ∖ sphere(8)
⊢ ground(peça)
\end{csg}

\csgfig{pre=flange-mini&size=900x520&dist=0.85}{flange}{0.9}{O programa acima: uma base com oito furos, quatro pilares e um rebaixo esférico.}

Os ingredientes já aparecem aí:
\begin{itemize}
  \item \textbf{definições} com \code{≔}, de constantes (\code{r}) e de funções (\code{furo(θ)});
  \item \textbf{isometrias} como valores: \code{Rz(θ)} é uma rotação, \code{T(…)} uma translação, compostas por \code{∘} e aplicadas por \code{·};
  \item \textbf{compreensão} no estilo de construção de conjuntos, \code{{ furo(k · 45°) | k ∈ 0..<8 }}, e a \textbf{união de uma família} \code{⋃};
  \item \textbf{operações de conjunto} entre sólidos: \code{∖} e \code{∪};
  \item o \textbf{objeto exibido}, marcado com \code{⊢}, que pode ser exportado como STL para impressão 3D.
\end{itemize}

A aplicação roda inteiramente no navegador. É escrita em HTML, CSS e JavaScript, sem etapa de compilação.
A única dependência é o núcleo geométrico \href{https://github.com/elalish/manifold}{manifold}, compilado para WebAssembly.
Ele garante que todo resultado seja um sólido fechado e bem orientado, pronto para o fatiador da impressora.

\begin{nota}
\textsf{\textbf{Convenções.}} As unidades são milímetros. O eixo $z$ aponta para cima e a mesa de impressão é o plano $z = 0$.
Ângulos são medidos em radianos: \code{45°} é apenas uma forma de escrever \code{45·π/180}.
\end{nota}

% =======================================================
\section{A aplicação}\label{sec:interface}

\subsection{Como abrir}
Por usar módulos JavaScript e WebAssembly, a aplicação precisa ser servida por HTTP; abrir o arquivo
\code{index.html} diretamente (\code{file://}) não funciona. Na pasta do projeto, execute
\code{python3 -m http.server} (ou \code{npm run serve}) e acesse \code{http://localhost:8000}.
A versão publicada no site funciona sem instalação.

\subsection{Visão geral da tela}
\begin{itemize}
  \item \textbf{Menu}, no topo:
    \begin{itemize}
      \item \emph{Arquivo}: novo, abrir e salvar programas \code{.csg}, copiar um link do programa, exportar STL;
      \item \emph{Editar}: desfazer, comentar linhas, substituições automáticas, paleta;
      \item \emph{Exibir}: grade, eixos, arestas, projeção, vistas, tema;
      \item \emph{Avaliar}: avaliação automática e qualidade;
      \item \emph{Exemplos};
      \item \emph{Ajuda}: referência rápida e este manual.
    \end{itemize}
  \item \textbf{Paleta de símbolos}, logo abaixo: um clique insere o símbolo no cursor. Passe o mouse para ver o nome, o atalho e a grafia ASCII.
  \item \textbf{Editor}, à esquerda: realce de sintaxe, números de linha, erros sublinhados e, abaixo dele, o painel de saída com erros e valores exibidos com \code{⊢}.
  \item \textbf{Visualização 3D}, à direita: o modelo sobre a grade da mesa de impressão, com o painel \textbf{Objetos} e a barra de vistas.
  \item \textbf{Barra de status}, embaixo: resultado e tempo da avaliação, número de triângulos, dimensões e volume do objeto selecionado, posição do cursor.
\end{itemize}

\subsection{Avaliação}
Com \emph{Avaliar ao digitar} ligado (padrão), o programa é reavaliado meio segundo depois da última tecla.
\kbd{Ctrl}+\kbd{Enter} avalia imediatamente. A avaliação acontece num processo separado (\emph{Web Worker}), então
a interface nunca trava. Se uma avaliação demorar mais de um segundo e você continuar editando, ela é cancelada.
Enquanto o programa tiver erro de sintaxe, a visualização mantém o último modelo válido.

\subsection{Visualização}
\begin{longtable}{@{}L{0.45}L{0.55}@{}}
\toprule
arrastar com o botão esquerdo & girar em torno do modelo \\
botão direito, botão do meio ou \kbd{Shift}+arrastar & deslocar \\
roda do mouse / pinça (toque) & aproximar e afastar \\
duplo clique, \kbd{F} ou \code{⤢} & enquadrar o modelo \\
Iso · Frente · Topo · Lado & vistas padrão \\
\bottomrule
\end{longtable}
Em \emph{Exibir} é possível ligar as arestas dos triângulos (útil para avaliar a resolução), desligar grade
e eixos, usar projeção ortográfica e escolher tema claro ou escuro.

\subsection{Painel de objetos}
Lista todos os \textbf{objetos nomeados} do programa (seção~\ref{sec:objetos}) e as expressões exibidas com \code{⊢}.
Para cada um há:
\begin{itemize}
  \item o botão \code{◉}/\code{○}, que exibe ou oculta o objeto;
  \item o botão \textbf{STL}, que exporta o objeto;
  \item a seleção, feita com um clique no nome, que define o objeto descrito na barra de status e exportado pelo menu \emph{Arquivo}.
\end{itemize}
Abaixo da lista aparecem os valores das constantes que não são sólidos.

\subsection{Arquivos e preferências}
O programa é salvo automaticamente no navegador; ao reabrir a página, ele volta. \kbd{Ctrl}+\kbd{S} baixa o
programa como arquivo \code{.csg} (texto UTF-8), que pode ser reaberto por \emph{Arquivo → Abrir}.
\emph{Arquivo → Copiar link do programa} copia um endereço que contém o programa inteiro, compactado:
quem abrir o link vê o mesmo programa, pronto para avaliar e modificar. Nada é enviado a um servidor; o
programa vai no próprio link (depois do \code{\#}). É um bom jeito de mandar um exercício ou uma solução.
As preferências de menus também são lembradas. Carregar um exemplo, abrir um link ou criar um programa novo pode ser desfeito com \kbd{Ctrl}+\kbd{Z}.

% =======================================================
\section{Digitando a notação matemática}\label{sec:simbolos}

Há quatro maneiras de escrever cada símbolo, e todas produzem exatamente o mesmo programa:
\begin{enumerate}
  \item \textbf{Atalho estilo LaTeX.} Digite \code{\cup} seguido de espaço, Tab ou qualquer sinal de pontuação e ele vira \code{∪}.
    Ao digitar \code{\}, aparece uma lista de sugestões que se filtra conforme você digita; navegue com as setas e confirme com
    \kbd{Tab} ou \kbd{Enter}. Letras gregas também funcionam: \code{\theta} → \code{θ},
    \code{\alpha} → \code{α}. Índices numéricos se digitam com \code{\_1} → \code{₁}.
  \item \textbf{Substituição automática.} Ao digitar \code{:=}, \code{->}, \code{<=}, \code{>=} ou \code{!=},
    o editor os troca por \code{≔ ↦ ≤ ≥ ≠}. Pode ser desligada em \emph{Editar}.
  \item \textbf{Paleta} de símbolos, com o mouse.
  \item \textbf{Grafia ASCII}, diretamente: a linguagem aceita \code{\/} para \code{∪}, \code{/\} para \code{∩},
    \code{\} para \code{∖}, \code{*} para \code{·}, \code{@} para \code{∘}, \code{and}, \code{or},
    \code{not}, \code{in}, \code{pi}, \code{45deg} etc. Os próprios nomes LaTeX (\code{\cup}) também são
    aceitos no código-fonte, mesmo sem a conversão.
\end{enumerate}

Os dois programas abaixo são idênticos para a linguagem:
\begin{csg}[run]
A ≔ cube(10) ∪ T(5, 0, 0) · sphere(6) ∖ cylinder(2, 30)
B := cube(10) \/ T(5, 0, 0) * sphere(6) \ cylinder(2, 30)
⊢ A, Tx(30) · B
\end{csg}

A tabela completa está na seção~\ref{sec:tabela}. Os mais usados:
\begin{longtable}{@{}L{0.3}L{0.70}@{}}
\toprule
\code{≔ ⊢ ↦} & \code{\def} ou \code{:=} · \code{\show} · \code{\to} ou \code{->} \\
\code{∪ ∩ ∖ △} & \code{\cup \cap \setminus \triangle} \\
\code{⋃ ⋂ Σ ∏} & \code{\bigcup \bigcap \sum \prod} \\
\code{∘ · × ⁻¹} & \code{\circ \cdot \times \inv} \\
\code{∈ ≤ ≥ ≠ ∧ ∨ ¬} & \code{\in \le \ge \ne \and \or \not} \\
\code{° π √ ² ⟨ ⟩} & \code{\deg \pi \sqrt \sq \langle \rangle} \\
\bottomrule
\end{longtable}

\subsection{Outros recursos do editor}
\begin{longtable}{@{}L{0.3}L{0.70}@{}}
\toprule
\kbd{Ctrl}+\kbd{Enter} & avaliar agora \\
\kbd{Ctrl}+\kbd{S} & salvar o programa \\
\kbd{Ctrl}+\kbd{/} & comentar ou descomentar as linhas selecionadas \\
\kbd{Tab} / \kbd{Shift}+\kbd{Tab} & indentar / desindentar (com várias linhas selecionadas) \\
\kbd{Enter} & nova linha mantendo a indentação \\
\kbd{F1} & referência rápida \\
\bottomrule
\end{longtable}
Erros aparecem sublinhados em vermelho no ponto exato, com o número da linha destacado. Passe o mouse sobre o número
para ver a mensagem, ou clique na mensagem no painel de saída para levar o cursor até ela.

% =======================================================
\section{Programas, definições e objetos}\label{sec:programas}

\subsection{Definições}
Um programa é uma sequência de \textbf{definições}, uma por linha:
\begin{csg}
nome ≔ expressão                        # constante
f(x, y) ≔ expressão                     # função
g(x, n = 3) ≔ expressão                 # parâmetro com valor padrão
\end{csg}
O sinal \code{=} também é aceito no lugar de \code{≔} numa definição. Fora das definições, \code{=} é comparação.
Nomes começam com letra (qualquer alfabeto, incluindo grego) ou \code{_} e podem conter dígitos, índices
(\code{r₁}), \code{_} e apóstrofos (\code{f′}). Comentários começam com \code{#} e vão até o fim da linha.

\subsection{A ordem não importa}
As definições de topo formam um sistema de equações, não uma sequência de comandos. Um nome pode ser usado antes
da linha em que é definido, e funções podem ser mutuamente recursivas. Cada constante é calculada uma única vez, quando
alguém precisa dela. Definições circulares (\code{a ≔ b}, \code{b ≔ a}) são detectadas e
geram erro. Um nome não pode ser definido duas vezes no mesmo programa. Já redefinir um nome embutido, como chamar um sólido
de \code{S}, é permitido, mas esconde a função original.
\begin{csg}[run]
peça ≔ base ∖ furo          # usa nomes definidos abaixo
base ≔ cylinder(raio, 5)
furo ≔ cylinder(raio / 3, 20)
raio ≔ 12
\end{csg}

\subsection{Objetos nomeados e o que é exibido}\label{sec:objetos}
Toda constante de topo cujo valor é um sólido é um \textbf{objeto nomeado}: aparece no painel de objetos, pode ser exibida e
pode ser exportada como STL com o seu nome. O comando \code{⊢} (ou \code{show}) escolhe o que a visualização mostra
inicialmente. Ele aceita várias expressões separadas por vírgula e pode aparecer em várias linhas.
\begin{csg}[run]
caixa  ≔ box(30, 20, 10)
tampa  ≔ Tz(15) · box(30, 20, 2)
⊢ caixa, tampa              # dois objetos, em cores diferentes
⊢ volume(caixa)             # números e listas aparecem no painel de saída
\end{csg}
Sem nenhum \code{⊢}, aparece o último objeto nomeado. Uma linha com uma expressão solta equivale a \code{⊢} dessa expressão.

\subsection{Linhas longas}\label{sec:linhas}
Uma definição pode ocupar várias linhas. Uma quebra de linha \textbf{não} termina a definição quando:
\begin{itemize}
  \item está dentro de parênteses, colchetes ou chaves ainda abertos;
  \item a linha seguinte começa indentada (com espaços);
  \item a linha termina com um operador, uma vírgula, \code{≔}, \code{↦}, \code{then}, \code{else}, \code{let}, \code{in};
  \item a linha seguinte começa com um operador binário (\code{∪}, \code{∖}, \code{·}, \code{|>}, …) ou com \code{then}/\code{else}/\code{in}.
\end{itemize}
Para recuperar melhor de erros, uma linha \textbf{não indentada} que começa com \code{nome ≔} sempre inicia uma nova definição,
mesmo depois de um parêntese esquecido. Várias definições curtas podem ficar na mesma linha separadas por \code{;}.
\begin{csg}[run]
roda(r, e) ≔
  let aro ≔ cylinder(r, e) ∖ cylinder(r - 2, 2e),
      raios ≔ ⋃{ Rz(k · 60°) · T(r/2, 0, 0) · box(r, 1.5, e) | k ∈ 0..<3 }
  in aro ∪ raios
     ∪ cylinder(3, e)
a ≔ 1; b ≔ 2
⊢ roda(20, 4)
\end{csg}

% =======================================================
\section{Valores e expressões}\label{sec:valores}

\subsection{Tipos de valores}
\begin{longtable}{@{}L{0.22}L{0.78}@{}}
\toprule
\colhead{Tipo} & \colhead{Exemplos} \\
\midrule
número & \code{3}, \code{-2.5}, \code{1e-3}, \code{π}, \code{τ} (= 2π), \code{e}, \code{∞}, \code{45°} \\
booleano & \code{true}, \code{false}, resultados de comparações \\
lista (e vetor) & \code{[1, 2, 3]}, \code{⟨1, 2, 3⟩}, \code{{-1, 1}}, \code{[]} \\
intervalo & \code{0..10}, \code{0..<10}, \code{range(0, 1, 0.25)} \\
sólido & \code{cube(10)}, \code{A ∪ B}, \code{∅} \\
perfil (região plana) & \code{circle(5)}, \code{square(4) ∖ circle(1)}, \code{polygon([⟨0, 0⟩, ⟨4, 0⟩, ⟨0, 3⟩])} \\
transformação & \code{T(1, 2, 3)}, \code{Rz(30°)}, \code{S(2)}, \code{id} \\
função & \code{x ↦ x²}, \code{sin}, \code{(∪)}, \code{cylinder(5)} (aplicação parcial) \\
\bottomrule
\end{longtable}
A tipagem é dinâmica. Uma operação aplicada a valores incompatíveis gera um erro que diz quais tipos foram encontrados,
por exemplo \emph{'∪' não se aplica a sólido e número}.

\subsection{Números}\label{sec:numeros}
\begin{itemize}
  \item \textbf{Aritmética}: \code{+ - · / ^}, com \code{*} e \code{×} como alternativas para o produto. \code{mod} é o resto matemático, sempre não negativo para divisor positivo: \code{-7 mod 3 = 2}.
  \item \textbf{Expoentes sobrescritos}: \code{x²}, \code{x³}, \code{2¹⁰}. A potência associa à direita: \code{2^3^2 = 2^9}. O sinal de menos tem precedência menor, então \code{-2² = -4}.
  \item \textbf{Raiz e arredondamentos}: \code{√x} ou \code{√(x + 1)}; \code{⌊x⌋} e \code{⌈x⌉} para piso e teto.
  \item \textbf{Graus}: \code{x°} multiplica por $\pi/180$. Funciona depois de qualquer expressão: \code{(a + b)°}.
  \item \textbf{Multiplicação implícita}: um número literal colado (sem espaço) a um nome, parêntese ou \code{√} multiplica:
    \code{2π}, \code{3x²} (= \code{3·x²}), \code{30°k}, \code{2(x + 1)}, \code{2√2}. Com espaço, \code{2 π} é erro.
    Entre dois nomes não há multiplicação implícita: escreva \code{a·b}.
  \item \textbf{Divisão por zero} e operações sem resultado real, como \code{√(-1)}, são erros, não \code{NaN}.
\end{itemize}
\begin{csg}[run]
⊢ 2π = τ, 3² + 4² = 5², 180° = π, ⌊7/2⌋, -7 mod 3, 1/2π
\end{csg}

\subsection{Listas, vetores e intervalos}
Colchetes, chaves e ângulos formam listas: \code{[a, b]}, \code{{a, b}} e \code{⟨a, b⟩} são o mesmo valor.
Use a grafia que comunica melhor a intenção: \code{⟨x, y, z⟩} para pontos e vetores, \code{{-1, 1}} para conjuntos de índices.
Uma lista de três números é um \textbf{vetor}:
\begin{itemize}
  \item \code{u + v}, \code{u - v} e \code{k · v} operam componente a componente;
  \item \code{u · v} é o produto escalar e \code{u × v} o vetorial;
  \item \code{v.x}, \code{v.y}, \code{v.z} e \code{v[0]} acessam componentes (índices começam em 0);
  \item \code{norm(v)}, \code{unit(v)} e \code{dist(p, q)} calculam norma, versor e distância.
\end{itemize}
\code{xs ++ ys} concatena listas. \code{∪}, \code{∩} e \code{∖} entre listas fazem a operação de conjuntos correspondente.

\textbf{Intervalos} são sequências preguiçosas de números:
\begin{itemize}
  \item \code{a..b} vai de $a$ até $b$ inclusive, com passo 1;
  \item \code{a..<b} exclui $b$;
  \item \code{range(a, b, passo)} aceita passo fracionário ou negativo.
\end{itemize}
Intervalos funcionam onde uma lista funciona: em compreensões, \code{map}, \code{∈}, \code{len}, \code{Σ}.
\begin{csg}[run]
v ≔ ⟨1, 2, 3⟩
⊢ v + [1, 1, 1], 2v, v · v, v × ⟨0, 0, 1⟩, v.z, len(0..<10), 3 ∈ 1..5, {1, 2} ∪ {2, 3}
\end{csg}

\subsection{Comparações e lógica}
\code{= ≠ < ≤ > ≥} comparam números; \code{=} e \code{≠} comparam qualquer valor estruturalmente (listas
elemento a elemento, transformações pela matriz). \code{x ∈ xs} e \code{x ∉ xs} testam pertinência em listas e intervalos.
Comparações podem ser \textbf{encadeadas} como na matemática: \code{0 ≤ i < n} significa \code{0 ≤ i ∧ i < n}.
Os conectivos são \code{¬ ∧ ∨} (ou \code{not and or}), com avaliação em curto-circuito. Condições precisam ser booleanas;
números não valem como verdadeiro ou falso.

\subsection{Precedência dos operadores}\label{sec:precedencia}
Do que liga mais forte para o que liga mais fraco:
\begin{longtable}{@{}L{0.45}L{0.55}@{}}
\toprule
\colhead{Operadores} & \colhead{Observação} \\
\midrule
\code{f(x)}\quad\code{xs[i]}\quad\code{v.x}\quad\code{x°}\quad\code{x²}\quad\code{g⁻¹} & pós-fixos \\
multiplicação implícita \code{2x} & só depois de número literal \\
\code{^} & associa à direita \\
\code{∘} & composição \\
\code{-x}\quad\code{√x} & prefixos \\
\code{·}\quad\code{×}\quad\code{/}\quad\code{mod} & produtos; \code{g · A} aplica transformação \\
\code{⋃}\quad\code{⋂}\quad\code{Σ}\quad\code{∏} & o operando inclui produtos, mas não somas \\
\code{+}\quad\code{-}\quad\code{++} & \\
\code{..}\quad\code{..<} & intervalos \\
\code{⊕} & soma de Minkowski \\
\code{∩} & como o produto em relação à soma \\
\code{∪}\quad\code{∖}\quad\code{△} & mesmo nível, associam à esquerda (como \code{+} e \code{-}) \\
\code{=}\quad\code{≠}\quad\code{<}\quad\code{≤}\quad\code{>}\quad\code{≥}\quad\code{∈}\quad\code{∉} & encadeáveis \\
\code{¬}, depois \code{∧}, depois \code{∨} & \\
\code{|>} & pipe \\
\code{↦}\quad\code{if … then … else}\quad\code{let … in} & estendem-se o mais à direita possível \\
\bottomrule
\end{longtable}
Consequências práticas:
\begin{itemize}
  \item \code{T(…) · A ∪ B} é \code{(T(…) · A) ∪ B}: a transformação só se aplica a \code{A}.
  \item \code{A ∪ B ∖ C} é \code{(A ∪ B) ∖ C}, porque \code{∪} e \code{∖} associam da esquerda para a direita.
  \item \code{A ∪ B ∩ C} é \code{A ∪ (B ∩ C)}.
  \item \code{Rz(θ) ∘ T(…) · A} compõe primeiro e aplica depois.
\end{itemize}
Na dúvida, use parênteses.

% =======================================================
\section{Sólidos e operações de conjunto}\label{sec:solidos}

\subsection{Primitivas}
Todas as primitivas são \textbf{centradas na origem}; use \code{ground} (seção~\ref{sec:consultas}) para apoiá-las na mesa. Os argumentos
podem ser posicionais ou nomeados: \code{cylinder(5, 20)} é o mesmo que \code{cylinder(h = 20, r = 5)}.
\begin{longtable}{@{}L{0.4}L{0.60}@{}}
\toprule
\colhead{Primitiva} & \colhead{Sólido} \\
\midrule
\code{cube(a)} & cubo de aresta $a$ (também \code{cube([x, y, z])}) \\
\code{box(x, y, z)} ou \code{box(v)} & paralelepípedo com essas dimensões \\
\code{sphere(r, n?)} & esfera (geodésica) de raio $r$ \\
\code{cylinder(r, h, n?)} & cilindro ao longo de $z$, de $-h/2$ a $h/2$ \\
\code{cone(r1, r2, h, n?)} & tronco de cone, raio $r_1$ embaixo e $r_2$ em cima (um deles pode ser 0) \\
\code{torus(R, r, n?, m?)} & toro em torno de $z$, raio maior $R$, raio do tubo $r < R$ \\
\code{prism(n, r, h)} & prisma regular de $n$ lados inscrito no círculo de raio $r$ \\
\code{polyhedron(pontos, faces)} & poliedro arbitrário (abaixo) \\
\code{hull(A, B, …)} & fecho convexo de sólidos e/ou pontos \\
\code{∅} (ou \code{empty}) & o sólido vazio, elemento neutro de $\cup$ \\
\bottomrule
\end{longtable}
O parâmetro opcional \code{n} fixa o número de segmentos usados para aproximar círculos. Sem ele, a resolução
segue a qualidade escolhida no menu \emph{Avaliar}: no máximo um certo ângulo por segmento e no mínimo um certo comprimento.
\code{prism(6, r, h)} é exatamente um prisma hexagonal, útil para encaixes de porcas.

\begin{csg}[run, id=prims]
⊢ ground(cube(10)), Tx(16) · ground(sphere(6)), Tx(31) · ground(cylinder(5, 12)),
  Tx(46) · ground(cone(6, 0, 12)), Tx(64) · ground(torus(7, 2.5)), Tx(82) · ground(prism(6, 6, 10)),
  Tx(98) · ground(hull(cube(6), T(4, 4, 8) · sphere(2)))
\end{csg}
\csgfig{pre=prims&size=1000x380&dist=0.75&view=iso}{primitivas}{1}{Cubo, esfera, cilindro, cone, toro, prisma hexagonal e um fecho convexo.}

\subsubsection{Poliedros}
\code{polyhedron} recebe a lista de vértices e a lista de faces. Cada face é uma lista de índices dos vértices (a partir de 0)
em \textbf{ordem anti-horária vista de fora}. A malha precisa fechar um volume: cada aresta pertence a exatamente duas faces.
Caso contrário, um erro é reportado.
\begin{csg}[run]
tetraedro ≔ polyhedron([[0, 0, 0], [10, 0, 0], [0, 10, 0], [0, 0, 10]],
                       [[0, 2, 1], [0, 1, 3], [0, 3, 2], [1, 2, 3]])
⊢ volume(tetraedro)       # 1000/6
\end{csg}

\subsection{Operações de conjunto}
\begin{longtable}{@{}L{0.32}L{0.68}@{}}
\toprule
\colhead{Operação} & \colhead{Significado} \\
\midrule
\code{A ∪ B} (também \code{A + B}) & união: pontos de $A$ ou de $B$ \\
\code{A ∩ B} & interseção: pontos de $A$ e de $B$ \\
\code{A ∖ B} (também \code{A - B}) & diferença: pontos de $A$ fora de $B$ \\
\code{A △ B} & diferença simétrica: $(A \setminus B) \cup (B \setminus A)$ \\
\code{A ⊕ B} & soma de Minkowski: $\{\, a + b \mid a \in A,\ b \in B \,\}$ (arredonda cantos) \\
\bottomrule
\end{longtable}
\begin{csg}[run, id=setops]
A ≔ cube(10)
B ≔ T(4, 4, 4) · sphere(6.5)
⊢ A ∪ B, Tx(25) · (A ∩ B), Tx(50) · (A ∖ B), Tx(75) · (A △ B)
\end{csg}
\csgfig{pre=setops&size=1000x360&dist=0.72}{conjuntos}{1}{Da esquerda para a direita: $A \cup B$, $A \cap B$, $A \setminus B$ e $A \mathbin{\triangle} B$, com $A$ um cubo e $B$ uma esfera.}
Os resultados são sólidos de verdade, regularizados: superfícies que só se tocam não deixam lâminas de espessura zero, e o
resultado é sempre uma malha fechada e orientada. Sólidos que apenas encostam numa face se fundem na união.

\subsection{Famílias: \texttt{⋃} e \texttt{⋂}}
\code{⋃ xs} une todos os sólidos de uma lista e \code{⋂ xs} intersecta todos. Listas aninhadas são achatadas.
Combinadas com compreensões (seção~\ref{sec:compreensao}), escrevem padrões inteiros numa linha. A união de uma família é calculada em lote pelo
núcleo, o que é muito mais rápido do que reduzir com \code{∪} um a um. \code{⋃} de uma lista vazia é \code{∅};
\code{⋂} de uma lista vazia é erro, pois seria o espaço inteiro. Os nomes \code{union(xs)} e
\code{intersection(xs)} fazem o mesmo, para uso em \emph{pipes}.

\subsection{Consultas e posicionamento}\label{sec:consultas}
\begin{longtable}{@{}L{0.3}L{0.70}@{}}
\toprule
\code{bounds(A)} & caixa envolvente \code{[[xmin, ymin, zmin], [xmax, ymax, zmax]]} \\
\code{size(A)} & dimensões \code{[dx, dy, dz]} \\
\code{volume(A)}, \code{area(A)} & volume (mm³) e área da superfície (mm²) \\
\code{isempty(A)} & se $A$ é o conjunto vazio \\
\code{center(A)} & translada $A$ para que o centro da caixa envolvente fique na origem \\
\code{ground(A)} & translada $A$ verticalmente para apoiá-lo em $z = 0$ \\
\bottomrule
\end{longtable}
Como os valores são calculados de verdade, eles podem dirigir o próprio projeto:
\begin{csg}[run]
corpo ≔ sphere(10) ∪ Tz(8) · cylinder(4, 10)
tampa ≔ Tz(size(corpo).z / 2 + 1) · cylinder(6, 2)    # encaixa acima do corpo
⊢ ground(corpo ∪ tampa)
⊢ bounds(corpo), volume(corpo)
\end{csg}

% =======================================================
\section{Isometrias, escalas e transformações afins}\label{sec:transf}

Transformações são valores como quaisquer outros: podem ser guardadas em constantes, passadas a funções, postas em listas
e comparadas. Internamente são matrizes afins $4 \times 4$.

\subsection{Isometrias}
\begin{longtable}{@{}L{0.45}L{0.55}@{}}
\toprule
\code{T(x, y, z)}, \code{T(v)}, \code{Tx(d)}, \code{Ty(d)}, \code{Tz(d)} & translações \\
\code{Rx(θ)}, \code{Ry(θ)}, \code{Rz(θ)} & rotações em torno dos eixos (regra da mão direita) \\
\code{R(eixo, θ)}, \code{R(eixo, θ, p)} & rotação em torno de um eixo qualquer, passando pela origem ou pelo ponto $p$ \\
\code{Mx}, \code{My}, \code{Mz} & reflexões nos planos $x = 0$, $y = 0$, $z = 0$ \\
\code{M(n)}, \code{M(n, p)} & reflexão no plano de normal $n$, pela origem ou pelo ponto $p$ \\
\code{id} & a identidade \\
\bottomrule
\end{longtable}
Rotações por múltiplos de 90° são calculadas de forma exata, sem erro de arredondamento.

\subsection{A álgebra das transformações}
\begin{longtable}{@{}L{0.25}L{0.75}@{}}
\toprule
\code{g · A} & aplica $g$ ao sólido $A$ (também a um ponto \code{⟨x, y, z⟩} ou a uma lista, elemento a elemento) \\
\code{g(A)} & o mesmo que \code{g · A} \\
\code{g ∘ h} & composição: primeiro $h$, depois $g$; \code{(g ∘ h) · A = g · (h · A)} \\
\code{g⁻¹}, \code{inverse(g)} & inversa \\
\code{g^n}, \code{gⁿ} & potência inteira (negativa usa a inversa); \code{g⁰ = id} \\
\code{g = h} & igualdade das transformações \\
\bottomrule
\end{longtable}
Como em \code{Rz(θ) ∘ T(d, 0, 0) · A}, a transformação mais próxima do sólido age primeiro. Aqui, o sólido é deslocado e depois girado em torno da origem.
Isometrias formam um grupo, e a linguagem permite raciocinar com ele:
\begin{csg}[run]
g ≔ Rz(90°) ∘ T(10, 0, 0)       # rotação de 90° em torno do ponto (5, 5, 0)
⊢ g⁴ = id, g · ⟨5, 5, 0⟩, (g ∘ g⁻¹) = id, Mx ∘ Mx = id
\end{csg}

\subsection{Escalas}
\begin{longtable}{@{}L{0.3}L{0.70}@{}}
\toprule
\code{S(k)} & escala uniforme a partir da origem (semelhança) \\
\code{S(x, y, z)}, \code{S(v)} & escala por eixo \\
\code{Sx(k)}, \code{Sy(k)}, \code{Sz(k)} & escala num só eixo \\
\code{scale(…)} & sinônimo de \code{S} \\
\code{at(g, p)} & $g$ em torno do ponto $p$: \code{T(p) ∘ g ∘ T(-p)} \\
\code{isometry(g)} & verdadeiro se $g$ preserva distâncias \\
\bottomrule
\end{longtable}
Escalas não são isometrias, mas pertencem ao mesmo tipo de valor e se compõem com elas pelas mesmas regras. A composição de
isometrias continua sendo isometria. Se uma escala entra na composição, o resultado é uma \textbf{transformação afim}, e as
mensagens de erro passam a usar esse nome. Um fator de escala zero é erro, pois achataria o sólido. Um fator negativo é uma
reflexão: \code{Sx(-1) = Mx}. A orientação das faces é corrigida automaticamente sempre que a transformação inverte o espaço.
\begin{csg}[run, id=affine]
ovo     ≔ S(1, 1, 1.4) · sphere(10)
torre   ≔ ⋃{ Tz(4k) · Rz(9°k) · S(1 - k/20, 1 - k/20, 1) · box(20, 20, 4.01) | k ∈ 0..<12 }
selo    ≔ cylinder(12, 3) ∖ at(S(0.6), [0, 0, 1.5]) · cylinder(12, 3)
⊢ ground(ovo), Tx(35) · ground(torre), Tx(70) · ground(selo)
⊢ isometry(Rz(9°) ∘ T(1, 2, 3)), isometry(Rz(9°) ∘ S(0.95))
\end{csg}
\csgfig{pre=affine&size=1000x420&dist=0.8}{afins}{0.95}{Um elipsoide (esfera escalada em $z$), uma torre de camadas giradas e encolhidas em $x$ e $y$, e um selo com rebaixo feito por escala em torno de um ponto.}

\subsection{Simetria}
Como transformações são valores, um \textbf{grupo de simetria} é simplesmente uma lista. A união da órbita de um motivo produz um
objeto com aquela simetria:
\begin{csg}[run, id=d4]
D₄ ≔ { Rz(k · 90°) ∘ s | k ∈ 0..3, s ∈ [id, My] }      # grupo diedral: 8 elementos
motivo ≔ T(14, 5, 0) · Rz(25°) · box(12, 3, 4) ∪ T(9, 9, 0) · sphere(2.5)
roseta ≔ ⋃{ g · motivo | g ∈ D₄ } ∪ cylinder(6, 4) ∖ cylinder(2, 10)
⊢ ground(roseta)
\end{csg}
\csgfig{pre=d4&size=800x560&view=top&dist=0.9}{d4}{0.6}{A órbita de um motivo sob o grupo diedral $D_4$, vista de cima.}

% =======================================================
\section{Funções e construções funcionais}\label{sec:funcional}

\subsection{Definindo e chamando funções}
\begin{csg}[run]
arruela(d, D = 2d, e = 1.5) ≔ cylinder(D/2, e) ∖ cylinder(d/2, 2e)
⊢ arruela(5), Tx(20) · arruela(5, e = 3), Tx(40) · arruela(d = 4, D = 15)
\end{csg}
\begin{itemize}
  \item Parâmetros com valor padrão vêm depois dos obrigatórios. O padrão pode usar os parâmetros anteriores (\code{D = 2d}).
  \item Argumentos nomeados (\code{e = 3}) vêm depois dos posicionais e podem estar em qualquer ordem.
  \item Funções podem ser recursivas e mutuamente recursivas. A profundidade é limitada para evitar laços infinitos.
  \item \code{f{…}} é atalho para chamar \code{f} com uma lista literal: \code{max{ i² | i ∈ 1..5 }}.
\end{itemize}

\subsection{Funções anônimas}
\code{x ↦ expressão} (ou \code{x -> expressão}) é uma função anônima. Com vários parâmetros, use parênteses:
\code{(x, y) ↦ x + y}. Um parâmetro pode ser desestruturado: \code{[a, b] ↦ a·b} recebe um par.
Funções anônimas capturam as variáveis visíveis onde foram escritas.

\subsection{Aplicação parcial}
Chamar uma função com \textbf{menos} argumentos posicionais do que ela exige devolve uma função que espera os restantes:
\code{cylinder(5)} é a função \code{h ↦ cylinder(5, h)}. Por isso as funções de listas recebem a lista por
último (\code{map(f, xs)}, \code{filter(p, xs)}): \code{map(f)} é uma função de listas, pronta para \emph{pipes}.
Se uma aplicação parcial chegar onde se esperava um sólido, a mensagem de erro diz qual argumento faltou.

\subsection{Operadores como funções: seções}
\begin{longtable}{@{}L{0.4}L{0.60}@{}}
\toprule
\code{(∪)}, \code{(+)}, \code{(·)}, \code{(=)}… & o operador binário como função de dois argumentos \\
\code{(A ∪)}, \code{(Rz(90°) ·)}, \code{(10 -)} & seção à esquerda: falta o operando da direita \\
\code{(∖ furo)}, \code{(· 2)}, \code{(mod 3)} & seção à direita: falta o operando da esquerda \\
\code{(⋃)}, \code{(Σ)}, \code{(√)}, \code{(¬)} & operadores prefixos como funções de um argumento \\
\bottomrule
\end{longtable}
\code{(- x)} é a negação de $x$, não uma seção.

\subsection{Composição e pipe}
\code{f ∘ g} compõe funções ($(f \circ g)(x) = f(g(x))$), assim como compõe transformações.
\code{x |> f} aplica $f$ a $x$. Encadeado, o pipe lê como uma linha de montagem:
\begin{csg}[run]
dente ≔ T(12, 0, 0) · box(4, 3, 5)
coroa ≔ 0..<12 |> map(k ↦ Rz(30°k) · dente) |> union
⊢ coroa ∪ cylinder(12, 5)
⊢ 1..10 |> filter(x ↦ x mod 2 = 0) |> map((· 10)), ((x ↦ x + 1) ∘ (x ↦ 2x))(5)
\end{csg}

\subsection{\texttt{map}, \texttt{filter}, \texttt{fold} e \texttt{reduce}}
\begin{longtable}{@{}L{0.25}L{0.75}@{}}
\toprule
\code{map(f, xs)} & \code{[f(x₀), f(x₁), …]} \\
\code{filter(p, xs)} & os elementos para os quais \code{p(x)} é verdadeiro \\
\code{fold(f, a, xs)} & \code{f(…f(f(a, x₀), x₁)…, xₙ)}, redução pela esquerda a partir de $a$ \\
\code{reduce(f, xs)} & como \code{fold}, começando do primeiro elemento (erro se vazia) \\
\code{iterate(f, k, x)} & \code{[x, f(x), f(f(x)), …]} com $k$ aplicações \\
\bottomrule
\end{longtable}
\begin{csg}[run]
peças ≔ [cube(8), T(6, 0, 0) · sphere(5), T(12, 0, 0) · cylinder(3, 12)]
⊢ reduce((∪), peças)
⊢ fold((+), 0, 1..100), iterate(x ↦ 2x, 5, 1)
\end{csg}

\subsection{\texttt{let} e \texttt{if}}
\code{let a ≔ …, b ≔ …, f(x) ≔ … in expressão} introduz nomes locais; cada definição enxerga as anteriores, e funções
locais podem ser recursivas. \code{if c then a else b} é uma expressão (o \code{else} é obrigatório); a condição
precisa ser booleana.
\begin{csg}[run, id=menger]
menger(k, s) ≔
  if k = 0 then cube(s)
  else ⋃{ T((s/3) · [i, j, l]) · menger(k - 1, s/3)
        | i, j, l ∈ -1..1, abs(i) + abs(j) + abs(l) > 1 }
⊢ ground(menger(2, 54))
\end{csg}
\csgfig{pre=menger&size=800x600&dist=0.85}{menger}{0.6}{A esponja de Menger de nível 2: recursão, compreensão em três índices e união de família.}

% =======================================================
\section{Compreensões e grandes operadores}\label{sec:compreensao}

\subsection{Compreensões}
Uma compreensão constrói uma lista à maneira da notação de conjuntos: \code{{ expressão | cláusulas }} ou
\code{[ expressão | cláusulas ]}. As cláusulas, separadas por vírgula, são processadas da esquerda para a direita:
\begin{longtable}{@{}L{0.32}L{0.68}@{}}
\toprule
\code{x ∈ A} (ou \code{x in A}) & \textbf{gerador}: $x$ percorre a lista ou intervalo $A$ \\
\code{i, j ∈ A} & dois geradores aninhados sobre o mesmo conjunto \\
\code{[a, b] ∈ pares} & gerador com desestruturação \\
\code{y ≔ expressão} & definição local, visível nas cláusulas seguintes e no corpo \\
qualquer condição booleana & \textbf{filtro}: mantém só as combinações que a satisfazem \\
\bottomrule
\end{longtable}
Geradores posteriores variam mais rápido, como laços aninhados. Um nome já introduzido por um gerador da mesma compreensão
faz de \code{x ∈ B} um filtro de pertinência, como na matemática: \code{{ x | x ∈ 1..20, x ∈ primos }}.
\begin{csg}[run]
⊢ { k² | k ∈ 1..10, k mod 2 = 1 }
⊢ [ [i, j] | i, j ∈ 0..2, i < j ]
⊢ { a + b | [a, b] ∈ [[1, 2], [3, 4]] }
⊢ { r | x ∈ 1..4, r ≔ √x, r = ⌊r⌋ }
\end{csg}

\begin{csg}[run, id=grid]
placa ≔ box(60, 60, 3)
furos ≔ ⋃{ T(10i, 10j, 0) · cylinder(1.5 + (i + j)/4, 10) | i, j ∈ -2..2, (i + j) mod 2 = 0 }
⊢ ground(placa ∖ furos)
\end{csg}
\csgfig{pre=grid&size=800x500&dist=0.85}{placa}{0.6}{Furos num padrão de tabuleiro, com diâmetro dependendo da posição.}

\subsection{Grandes operadores}
\code{⋃}, \code{⋂}, \code{Σ} e \code{∏} aplicam-se a uma coleção, ou a uma compreensão escrita no índice:
\begin{csg}
⋃ xs                 ⋃{ e | i ∈ A }           # sobre uma lista/compreensão
⋃_{i ∈ A, cond} e    Σ_{k ∈ 1..n} k²          # índice com cláusulas
Σ_{k=1}^{n} k²       ∏_{k=1}^{5} k            # forma clássica com limites
\end{csg}
\code{Σ} soma números ou vetores (\code{Σ} de vazio é 0); \code{∏} multiplica (de vazio é 1). O corpo de um grande
operador vai até o próximo \code{+}, \code{-}, \code{∪} ou operador mais fraco: \code{Σ_{k=1}^{3} 2k + 1} é
\code{(Σ 2k) + 1}, assim como na escrita matemática, em que $\sum_{k=1}^{3} 2k + 1 = 13$.
\begin{csg}[run]
⊢ Σ_{k=1}^{10} k², ∏_{k=1}^{5} k, Σ{ ⟨k, 0, 1⟩ | k ∈ 1..3 }
anéis ≔ ⋃_{k ∈ 1..4} torus(8k, 1.5)
⊢ anéis
\end{csg}

% =======================================================
\section{Perfis: extrusão e revolução}\label{sec:perfis}
Muitas peças são mais fáceis de desenhar no plano: o contorno de uma letra, o corte de um vaso, uma estrela. Um
\textbf{perfil} é uma região do plano $z = 0$. Ele não é impresso diretamente: vira sólido por \textbf{extrusão}
(empilhado ao longo de $z$) ou por \textbf{revolução} (girado em torno de um eixo).

\subsection{Construindo perfis}
Os pontos do plano se escrevem com duas coordenadas, \code{⟨x, y⟩}.
\begin{longtable}{@{}L{0.36}L{0.64}@{}}
\toprule
\colhead{Perfil} & \colhead{Região} \\
\midrule
\code{circle(r, n?)} & disco de raio $r$ centrado na origem \\
\code{square(a)} ou \code{square(x, y)} & quadrado ou retângulo centrado na origem \\
\code{polygon(pontos)} & região limitada pela poligonal fechada (a orientação não importa) \\
\code{polygon([c₁, c₂, …])} & vários contornos: os interiores viram furos (regra par–ímpar) \\
\code{region(f, a, b)} & região entre o gráfico de $f$ e o eixo $x$, para $a \le x \le b$ \\
\code{offset(P, d)} & engorda o perfil de $d$ (afina se $d < 0$), arredondando os cantos \\
\code{hull(P, Q, …)} & fecho convexo de perfis e/ou pontos \code{⟨x, y⟩} \\
\bottomrule
\end{longtable}
As operações de conjunto \code{∪}, \code{∩}, \code{∖}, \code{△}, os grandes operadores \code{⋃} e \code{⋂} e as
transformações que preservam o plano (\code{T(x, y)}, \code{Rz}, \code{Mx}, \code{My}, \code{S(k)}…) funcionam com perfis
exatamente como com sólidos. \code{area}, \code{size}, \code{bounds} e \code{center} também. Perfis e sólidos não se
misturam numa mesma operação: a mensagem de erro sugere extrudar ou revolucionar o perfil antes.

Um perfil exibido com \code{⊢} aparece como uma placa fina sobre a mesa, e a barra de status mostra a sua área.
\begin{csg}[run, id=perfis]
estrela(n, R, r) ≔ polygon([ ρ · ⟨cos(θ), sin(θ)⟩
                     | k ∈ 0..<2n, θ ≔ k · π/n, ρ ≔ if k mod 2 = 0 then R else r ])
anel ≔ circle(10) ∖ circle(6)
L ≔ polygon([⟨0, 0⟩, ⟨16, 0⟩, ⟨16, 5⟩, ⟨5, 5⟩, ⟨5, 20⟩, ⟨0, 20⟩])
⊢ estrela(5, 12, 5), Tx(30) · anel, T(48, -10) · L, Tx(80) · offset(estrela(5, 12, 5), 2)
⊢ area(anel)            # π(10² − 6²) ≈ 201
\end{csg}
\csgfig{pre=perfis&size=1000x340&dist=0.8&view=top}{perfis}{0.95}{Uma estrela, um anel, um perfil em L e a estrela arredondada por \texttt{offset}.}

\subsection{Extrusão}
\code{extrude(P, h)} empilha o perfil de $z = 0$ até $z = h$, formando um prisma generalizado. Dois argumentos
opcionais deformam o empilhamento: \code{θ} gira o topo em relação à base (a torção é uniforme ao longo da altura) e
\code{k} escala o topo (\code{k = 0} fecha numa ponta, como uma pirâmide; \code{⟨kx, ky⟩} escala cada eixo).
\begin{longtable}{@{}L{0.36}L{0.64}@{}}
\toprule
\code{extrude(P, h)} & prisma de altura $h$ \\
\code{extrude(P, h, θ)} & prisma torcido: o topo gira $θ$ \\
\code{extrude(P, h, θ, k)} & torcido e com o topo escalado por $k$ \\
\bottomrule
\end{longtable}
\begin{csg}[run, id=extrusoes]
estrela(n, R, r) ≔ polygon([ ρ · ⟨cos(θ), sin(θ)⟩
                     | k ∈ 0..<2n, θ ≔ k · π/n, ρ ≔ if k mod 2 = 0 then R else r ])
P ≔ estrela(5, 12, 6)
⊢ extrude(P, 8), Tx(30) · extrude(P, 25, 72°), Tx(60) · extrude(P, 25, 0°, 0),
  Tx(90) · extrude(square(16), 25, 90°, 0.4)
\end{csg}
\csgfig{pre=extrusoes&size=1000x400&dist=0.75}{extrusoes}{0.95}{Extrusão simples, torcida (72°), em ponta (\texttt{k = 0}) e torcida com o topo reduzido.}
A extrusão em ponta é uma pirâmide de base $P$; seu volume é $\frac13 \cdot \text{área}(P) \cdot h$, o que se confere com
\code{volume} e \code{area}. A torção não altera o volume: cada fatia horizontal continua sendo uma cópia girada de $P$
(princípio de Cavalieri).

\subsection{Revolução}
\code{revolve(P)} gira o perfil em torno do \textbf{eixo $y$} do plano, que passa a ser o eixo $z$ do sólido: um ponto
\code{⟨x, y⟩} do perfil descreve o círculo de raio $x$ na altura $y$. Por isso o perfil deve estar no semiplano
$x \ge 0$; a parte com $x < 0$ é descartada. \code{revolve(P, θ)} gira só o ângulo $θ$, o que mostra o corte do sólido.

Num vaso, o raio é função da altura, $r(z)$: o perfil é a região $0 \le x \le r(y)$, escrita como um polígono por
compreensão. A soma dos volumes de discos finos aproxima a integral $\pi\int r(z)^2\,dz$, e o resultado se compara com o
volume calculado pelo núcleo:
\begin{csg}[run, id=vaso]
H ≔ 80;  e ≔ 2
r(z) ≔ 20 + 7sin(z/10) + z/6       # raio na altura z
externo ≔ polygon([⟨0, 0⟩] ++ [⟨r(z), z⟩ | z ∈ range(0, H, 0.5)] ++ [⟨0, H⟩])
interno ≔ polygon([⟨0, e⟩] ++ [⟨r(z) - e, z⟩ | z ∈ range(e, H + 1, 0.5)] ++ [⟨0, H + 1⟩])
vaso ≔ revolve(externo ∖ interno)
⊢ vaso, Tx(100) · revolve(externo ∖ interno, 270°)
⊢ volume(revolve(externo)), Σ[ π · r(z)^2 · 0.5 | z ∈ range(0.25, H, 0.5) ]
\end{csg}
\csgfig{pre=vaso&size=1000x480&dist=0.8}{vaso-revolucao}{0.85}{O vaso e a mesma revolução feita só até 270°, que mostra a parede e o fundo.}

Com \code{region}, a região sob um gráfico girada em torno do eixo $y$ dá o \textbf{método das cascas}:
$V = 2\pi\int_a^b x\,f(x)\,dx$. Para $f(x) = x$ em $[0, 2]$, $V = 16\pi/3 \approx 16{,}76$:
\begin{csg}[run]
⊢ volume(revolve(region(x ↦ x, 0, 2), 360°, 400)), 16π/3
\end{csg}
Já o teorema de Pappus diz que girar uma região de área $A$ cujo centroide está à distância $d$ do eixo gera volume
$2\pi d A$: um quadrado de lado 2 centrado em $x = 10$ dá um anel de volume $80\pi$.
\begin{csg}[run]
⊢ volume(revolve(Tx(10) · square(2), 360°, 400)), 80π
\end{csg}

% =======================================================
\section{Exemplos comentados}\label{sec:exemplos}
Todos estão no menu \emph{Exemplos} da aplicação.

\subsection{Engrenagem}
\begin{csg}[run, id=gear]
z ≔ 18; R ≔ 20; e ≔ 6
dente  ≔ T(R, 0, 0) · hull(box(4, 3.4, e), T(3, 0, 0) · box(1, 1.4, e))
dentes ≔ { Rz(k · τ/z) · dente | k ∈ 0..<z }
alívios ≔ ⋃{ Rz(k · 60°) · T(12.5, 0, 0) · cylinder(3.5, 2e) | k ∈ 0..<6 }
engrenagem ≔ reduce((∪), [cylinder(R, e)] ++ dentes)
  ∖ alívios ∪ cylinder(r = 7, h = e + 4) ∖ prism(6, 3.2, 3e)
⊢ ground(engrenagem)
\end{csg}
O dente trapezoidal é o fecho convexo de duas caixas. A compreensão gera as $z$ cópias giradas por múltiplos de
$\tau/z$. Os alívios circulares são retirados da coroa, o cubo central é recolocado, e o furo hexagonal recebe uma porca ou um eixo sextavado.
\csgfig{pre=gear&size=800x520&dist=0.8}{engrenagem}{0.6}{Engrenagem de 18 dentes.}

\subsection{Vaso torcido}
\begin{csg}[run, id=vase]
n ≔ 36; h ≔ 2; parede ≔ 2
perfil(k) ≔ 18 + 5·sin(k/6)
camada(k, recuo) ≔ Tz(k·h) · Rz(k · 3°) · prism(6, perfil(k) - recuo, h + 0.01)
vaso ≔ ⋃{ camada(k, 0) | k ∈ 0..<n } ∖ ⋃{ camada(k, parede) | k ∈ 1..n }
⊢ ground(vaso)
\end{csg}
Cada camada é um prisma hexagonal girado 3° em relação à anterior, com raio dado por um perfil senoidal. As camadas se sobrepõem
0,01\,mm para que a união seja um único sólido. O interior é a mesma pilha com raio menor, deslocada uma camada para cima: assim o
fundo fica fechado e a boca aberta.
\csgfig{pre=vase&size=800x600&dist=0.85}{vaso}{0.55}{Vaso torcido de 36 camadas.}

\subsection{Caixa com tampa}
\begin{csg}[run]
caixa(l, w, h, e = 2) ≔
  let externa ≔ box(l, w, h),
      interna ≔ Tz(e) · box(l - 2e, w - 2e, h)
  in externa ∖ interna
folga ≔ 0.2
fundo ≔ ground(caixa(60, 40, 30))
tampa ≔ T(0, 50, 0) · ground(box(60, 40, 2)
        ∪ Tz(2) · box(60 - 4 - 2folga, 40 - 4 - 2folga, 4)
        ∖ Tz(2) · box(60 - 8, 40 - 8, 6))
⊢ fundo, tampa
\end{csg}
Uma função parametrizada com espessura padrão. A tampa tem uma aba que entra na caixa com \code{folga} de
0,2\,mm de cada lado. Os dois objetos nomeados são exportados separadamente (botão STL de cada um).

% =======================================================
\section{Do modelo à impressão 3D}\label{sec:impressao}

\subsection{Exportando STL}
Escolha o objeto no painel \textbf{Objetos} e clique em \textbf{STL} ao lado dele, ou use \emph{Arquivo → Exportar STL do objeto selecionado}.
O arquivo recebe o nome do objeto (\code{peça.stl}). Expressões sem nome, exibidas com \code{⊢}, são exportadas como
\code{modelo.stl}. Por padrão o formato é STL binário, compacto. Há também STL ASCII, legível, e a exportação de todos os
objetos exibidos de uma vez. Para exportar uma variante, dê um nome a ela: \code{final ≔ ground(peça)}.

\subsection{Garantias}
Todo sólido produzido pelo núcleo é uma \textbf{variedade fechada e orientada}: cada aresta pertence a exatamente dois triângulos,
e as normais apontam para fora. Não há furos nem faces invertidas, e o fatiador não precisa reparar a malha. As coordenadas estão em milímetros.

\subsection{Recomendações}
\begin{itemize}
  \item \textbf{Apoie na mesa}: as primitivas são centradas na origem; use \code{ground(peça)} antes de exportar.
  \item \textbf{Qualidade}: \emph{Avaliar → Qualidade} controla a resolução das curvas (rascunho, normal, fina). Use rascunho enquanto modela e
    fina para exportar. Para um círculo específico, fixe \code{n}: \code{cylinder(2, 10, n = 64)}.
  \item \textbf{Sobreposição}: ao subtrair um furo que atravessa a peça, faça a ferramenta ultrapassar as faces. Por exemplo,
    \code{cylinder(2, h + 2)} para uma placa de altura $h$. Ao unir peças empilhadas, sobreponha-as um pouco.
    O núcleo lida com faces coincidentes, mas a sobreposição deixa a intenção explícita.
  \item \textbf{Folgas}: peças que se encaixam precisam de folga, tipicamente 0,1 a 0,3\,mm conforme a impressora. Parametrize-a.
  \item \textbf{Espessuras mínimas}: paredes mais finas que duas larguras de extrusão (cerca de 0,8\,mm com bico de 0,4\,mm) tendem a falhar.
  \item \textbf{Balanços}: superfícies inclinadas mais de \textasciitilde45° em relação à vertical precisam de suporte. Gire a peça com \code{Rx}/\code{Ry} para uma orientação favorável.
  \item \textbf{Verifique}: a barra de status mostra as dimensões e o volume do objeto selecionado, e \code{⊢ size(peça)} imprime as medidas no painel de saída.
\end{itemize}

% =======================================================
\section{Erros, limites e desempenho}\label{sec:erros}

\subsection{Mensagens de erro}
Erros de sintaxe são detectados antes da avaliação. Todos os erros de sintaxe do programa são listados de uma vez, e a
visualização mantém o último modelo válido. Erros de avaliação são reportados por definição: uma constante com erro não impede
as outras de serem calculadas, e as definições que dependem dela não repetem a mensagem. Exemplos típicos:
\begin{longtable}{@{}L{0.48}L{0.52}@{}}
\toprule
\colhead{Mensagem} & \colhead{Causa provável} \\
\midrule
\emph{'x' não está definido} & nome digitado errado ou definição faltando \\
\emph{'∪' não se aplica a sólido e número} & operandos de tipos incompatíveis \\
\emph{… (aplicação parcial de cylinder, faltando h)} & faltou um argumento obrigatório \\
\emph{cylinder não tem parâmetro 'q'} & argumento nomeado inexistente; a mensagem lista os válidos \\
\emph{para interseção use ∩} & \code{A · B} entre dois sólidos \\
\emph{a isometria vem à esquerda: g · S} & \code{A · g} em vez de \code{g · A} \\
\emph{definição circular envolvendo 'a'} & \code{a} depende de si mesma \\
\emph{recursão profunda demais} & recursão sem caso base alcançável \\
\emph{esperava ')' …} & parêntese não fechado; veja a posição sublinhada \\
\emph{tuplas não existem…} & \code{(a, b)} fora de lambda; use \code{[a, b]} \\
\bottomrule
\end{longtable}

\subsection{Limites}
Para proteger contra laços infinitos, há um limite de passos de avaliação (cerca de 20 milhões), de profundidade de
recursão (400 chamadas aninhadas) e de tamanho de listas (um milhão de elementos). Além disso, uma avaliação longa é
cancelada quando o programa é editado.

\subsection{Desempenho}
O núcleo manifold é rápido: os exemplos deste manual são avaliados em poucas dezenas de milissegundos. Para modelos grandes:
\begin{itemize}
  \item prefira \code{⋃ lista} a encadear \code{∪} um a um, pois a união de família é feita em lote;
  \item use qualidade \emph{rascunho} durante a modelagem;
  \item desligue \emph{Avaliar ao digitar} e avalie com \kbd{Ctrl}+\kbd{Enter};
  \item funções como \code{volume} e \code{bounds} forçam o cálculo da geometria. Evite chamá-las dentro de laços grandes.
\end{itemize}

% =======================================================
\section{Referência das funções embutidas}\label{sec:referencia}
Gerada a partir do código da aplicação. Parâmetros marcados com \code{?} são opcionais. Todas as funções aceitam
argumentos nomeados e aplicação parcial (exceto as de número variável de argumentos, marcadas com \code{…}).
Constantes: \code{π pi τ tau e ∞ inf ∅ empty id Mx My Mz}.

\InputIfFileExists{gen/functions.tex}{}{\emph{(execute \texttt{npm run manual} para gerar esta referência)}}

% =======================================================
\section{Tabela de símbolos}\label{sec:tabela}
Gerada a partir do código da aplicação. Os atalhos funcionam no editor (seguidos de espaço, Tab ou pontuação) e também
diretamente no código-fonte. A grafia ASCII é aceita pela linguagem. As letras gregas têm atalhos com o nome em inglês:
\code{\alpha \beta \gamma \delta \epsilon \theta \lambda \mu \rho \sigma \phi \psi \omega} e maiúsculas
\code{\Gamma \Delta \Theta \Lambda \Phi \Psi \Omega}.

\InputIfFileExists{gen/symbols.tex}{}{\emph{(execute \texttt{npm run manual} para gerar esta tabela)}}

% =======================================================
\section{Gramática resumida}\label{sec:gramatica}
\begin{csg}[plain]
programa    ::= { instrução (fim-de-linha | ";") }
instrução   ::= nome ("≔" | "=") expr
              | nome "(" parâmetros ")" ("≔" | "=") expr
              | "⊢" expr { "," expr }
              | expr
parâmetros  ::= [ nome { "," nome } ] { "," nome "=" expr }

expr        ::= nome "↦" expr | "(" nomes ")" "↦" expr | "[" nomes "]" "↦" expr
              | "if" expr "then" expr "else" expr
              | "let" ligação { "," ligação } "in" expr
              | expr op-binário expr | op-prefixo expr | expr pós-fixo
              | grande-op expr | grande-op "_{" cláusulas "}" expr
              | grande-op "_{" nome "=" expr "}^{" expr "}" expr
              | número | nome | "true" | "false"
              | "(" expr ")" | seção
              | abre [ expr { "," expr } ] fecha
              | abre expr "|" cláusulas fecha
              | "⌊" expr "⌋" | "⌈" expr "⌉"
pós-fixo    ::= "(" argumentos ")" | "[" expr "]" | "." nome | "°" | sobrescrito | "{" … "}"
argumentos  ::= [ expr { "," expr } ] { "," nome "=" expr }
cláusulas   ::= cláusula { "," cláusula }
cláusula    ::= nome { "," nome } ("∈" | "in") expr | "[" nomes "]" ("∈" | "in") expr
              | nome "≔" expr | expr
seção       ::= "(" op ")" | "(" expr op ")" | "(" op expr ")"
abre/fecha  ::= "[" "]" | "{" "}" | "⟨" "⟩"
grande-op   ::= "⋃" | "⋂" | "Σ" | "∏"
op-prefixo  ::= "-" | "√" | "¬"
op-binário  ::= "|>" "∨" "∧" "=" "≠" "<" "≤" ">" "≥" "∈" "∉" "∪" "∖" "△" "∩" "⊕"
                ".." "..<" "+" "-" "++" "·" "×" "/" "mod" "∘" "^"
\end{csg}
A precedência dos operadores binários está na seção~\ref{sec:precedencia}. Comentários (\code{#} até o fim da linha) e espaços são ignorados,
exceto pelas regras de continuação de linha (seção~\ref{sec:linhas}) e de multiplicação implícita (seção~\ref{sec:numeros}).

\end{document}
